\scanpage{68}
c) \deleted{\unsure{Élim} Si on remplace ``$f$} \inserted{Élimination de l'hypothèse :} ``$f$ de présentation finie'', \deleted{``\unsure{à remplacer}} (il suffit $f$ quasi fini-séparé).

\deleted{Ou bien} En supposant $f$ quasi-projectif à priori, et faisant un argument de passage à la limite dans $X=\bigcap\limits_i X_i$.

\deleted{Ou bien on \unsure{prou|peut}} Mais on prouve par ailleurs que $f$ est nécessairement quasi-affine (donc quasi-projectif) par la technique de localisation étale\inserted{ (\uline{IV}, 18.12.12)}, en utilisant pour ``descendre'' la prop \uline{IV} 17.7.4. \inserted{i.e. \uline{IV},11.3.16} (assez délicate, \deleted{\illegible{1}} utilisant le passage à la limite de la platitude). // On notera que la technique de localisation étale \uline{utilise} cependant le Main Theorem (local en haut) tel qu'il est contenu dans b) ; à ce stade, on a le choix entre sa démonstration \unsure{chronologique|cohomologique} \inserted{globale}, ou sa démonstration purement locale ....
% archive sheet number: 67
